Los resultados esperados se organizan según las tres perspectivas solicitadas por el programa de Maestría en Inteligencia Artificial.

\section{Generación de nuevo conocimiento}

\section{Generación de Nuevo Conocimiento}

\begin{table}[htbp!]
    \centering
    \caption{Resultados o productos esperados de nuevo conocimiento.}
    \pgfplotstabletypeset[
	  col sep=colon, 
	  trim cells=true,
	  column type={|p{0.52\textwidth}|p{0.22\textwidth}|p{0.18\textwidth}}, 
	  columns/Resultado/.style = {column name={\cellcolor{AzulInstitucional}\textbf{\color{Blanco}Resultado o Producto esperado}}}, 
	  columns/Indicador/.style = {column name={\cellcolor{AzulInstitucional}\textbf{\color{Blanco}Indicador }}},
	  columns/Beneficiario/.style = {column name={\cellcolor{AzulInstitucional}\textbf{\color{Blanco}Beneficiario}}},
	  every head row = { \hline},
	  after row = {\hline},
	  string type,
	  skip rows between index={10}{20}
	]
	{Tables/NuevoConocimiento.csv}
    \label{tab:resEspNuecoConocimiento}
\end{table}


\section{Fortalecimiento de la capacidad científica nacional}

\begin{itemize}
    \item \textbf{Formación de recurso humano:} Formación de un magíster en Inteligencia Artificial con competencias en procesamiento de lenguaje natural, sistemas multi-agente y verificación automatizada de hechos.
    \item \textbf{Colaboración interdisciplinaria:} El proyecto integra conocimientos de inteligencia artificial, lingüística computacional, regulación financiera y periodismo de verificación, fortaleciendo la capacidad de investigación interdisciplinaria de la Universidad Sergio Arboleda.
    \item \textbf{Infraestructura de investigación:} El benchmark de evaluación y el pipeline de recolección automatizada quedan disponibles como recursos para futuras investigaciones del grupo de investigación.
\end{itemize}

\section{Apropiación social del conocimiento}

\begin{itemize}
    \item \textbf{Herramienta de apoyo a verificadores:} El sistema puede ser utilizado por organizaciones de fact-checking colombianas (ColombiaCheck, AFP Factual) como herramienta de priorización y pre-verificación de contenido financiero sospechoso.
    \item \textbf{Repositorio público:} El código fuente, el benchmark y la documentación se publicarán en un repositorio de GitHub bajo licencia abierta, permitiendo la replicación y extensión por parte de la comunidad.
    \item \textbf{Divulgación:} Presentación de resultados en la ExpoPosgrados ECEI y en al menos un evento académico nacional o internacional (candidatos: IberLEF, encuentro nacional de IA).
    \item \textbf{Impacto social:} Contribución a la protección del ciudadano colombiano frente a la desinformación financiera, particularmente en poblaciones con menor alfabetización financiera y mediática.
\end{itemize}
